% Author: Izaak Neutelings (October 2021)
% Inspiration
%   "Very special relativity - An illustrated guide", Sander Bais (2007)
%   http://people.uncw.edu/hermanr/GR/Minkowski/Minkowski.pdf
\documentclass[border=3pt,tikz]{standalone}
\usepackage{amsmath} % for \text
\usepackage{etoolbox} % ifthen
\usepackage[outline]{contour} % glow around text
\usetikzlibrary{calc} % for adding up coordinates
\usetikzlibrary{decorations.markings,decorations.pathmorphing}
\usetikzlibrary{angles,quotes} % for pic (angle labels)
\usetikzlibrary{arrows.meta} % for arrow size
\usepackage{xfp} % higher precision (16 digits?)
\contourlength{1.1pt}

\tikzset{>=latex} % for LaTeX arrow head
\colorlet{myred}{red!85!black}
\colorlet{mydarkred}{red!55!black}
\colorlet{mylightred}{red!85!black!12}
\colorlet{myfieldred}{mydarkred!5} % for S' background
\colorlet{myredhighlight}{myred!20} % highlights simultaneity in ladder paradox
\colorlet{myblue}{blue!80!black}
\colorlet{mydarkblue}{blue!50!black}
\colorlet{mylightblue}{blue!50!black!30}
\colorlet{mylightblue2}{myblue!10}
\colorlet{mygreen}{green!80!black}
\colorlet{mypurple}{blue!40!red!80!black}
\colorlet{mydarkgreen}{green!50!black}
\colorlet{mydarkpurple}{blue!40!red!50!black}
\colorlet{myorange}{orange!40!yellow!95!black}
\colorlet{mydarkorange}{orange!40!yellow!85!black}
\colorlet{mybrown}{brown!20!orange!90!black}
\colorlet{mydarkbrown}{brown!20!orange!55!black}
\colorlet{mypurplehighlight}{mydarkpurple!20} % highlights simultaneity in ladder paradox
\tikzstyle{world line}=[myblue!40,line width=0.3]
\tikzstyle{world line t}=[mypurple!50!myblue!40,line width=0.3]
\tikzstyle{world line'}=[mydarkred!40,line width=0.3]
\tikzstyle{mysmallarr}=[-{Latex[length=3,width=2]},thin]
\tikzstyle{mydashed}=[dash pattern=on 3 off 3]
\tikzstyle{rod}=[mydarkbrown,draw=mydarkbrown,double=mybrown,double distance=2pt,
                 line width=0.2,line cap=round,shorten >=1pt,shorten <=1pt]
%\tikzstyle{rod'}=[rod,draw=mydarkbrown!80!red!85,double=mybrown!80!red!85]
\tikzstyle{vector}=[->,line width=1,line cap=round]
\tikzstyle{vector'}=[vector,shorten >=1.2]
\tikzstyle{particle}=[mygreen,line width=0.9]
\tikzstyle{photon}=[-{Latex[length=5,width=4]},myorange,line width=0.8,decorate,
                    decoration={snake,amplitude=1.0,segment length=5,post length=5}]

\def\tick#1#2{\draw[thick] (#1) ++ (#2:0.06) --++ (#2-180:0.12)}
\def\tickp#1#2{\draw[thick,mydarkred] (#1) ++ (#2:0.06) --++ (#2-180:0.12)}
\def\Nsamples{100} % number samples in plot

\begin{document}

\def\axes{ % common axes
  
  % SETTINGS
  \def\xmin{0.3}
  \def\xmax{3.1}
  \def\ymin{0.3}
  \def\ymax{2.6}
  \def\xminp{0.4} % minimum of rotated axis
  \def\xmaxp{2.85} % maximum of rotated axis
  \def\Nlines{4} % number of world lines (at constant x/t)
  \pgfmathsetmacro\ang{atan(0.5)} % angle between x and x' axes
  \pgfmathsetmacro\d{0.64*\xmax/\Nlines} % grid size
  \pgfmathsetmacro\D{\d/cos(\ang)/sqrt(1-tan(\ang)^2)} % grid size, boosted
  \pgfmathsetmacro\dextra{(\Nlines+1)*\d} % extra line
  \pgfmathsetmacro\st{3*\d} % spacetime interval
  \pgfmathsetmacro\sx{4*\d} % spacetime interval
  \pgfmathsetmacro\sr{sqrt(\sx^2-\st^2)} % spacetime interval sr^2 = st^2 - sx^2 < 0
  \pgfmathsetmacro\Ax{3*\D*sin(\ang)} % x coordinate of event A
  \pgfmathsetmacro\Ay{3*\D*cos(\ang)} % y coordinate of event A
  \pgfmathsetmacro\Bx{4*\D*cos(\ang)} % x coordinate of event B
  \pgfmathsetmacro\By{4*\D*sin(\ang)} % y coordinate of event B
  \pgfmathsetmacro\Cx{\Ay+\By} % x coordinate of event C' %(\Bx+tan(\ang)*\Ay)/sqrt(1-tan(\ang)^2)
  \coordinate (O)  at (0,0);
  \coordinate (X)  at (\xmax+0.2,0);
  \coordinate (T)  at (0,\ymax+0.2);
  \coordinate (X') at (\ang:\xmaxp+0.2);
  \coordinate (T') at (90-\ang:\xmaxp+0.2);
  \coordinate (A)  at (0,\st);        % event A
  \coordinate (A') at (90-\ang:3*\D); % event A', boosted A
  \coordinate (B)  at (\sx,0);        % event A
  \coordinate (B') at (\ang:4*\D);    % event A', boosted A
  \coordinate (C)  at (4*\d,3*\d);    % event C
  
  % WORLD LINE GRID
  \message{  Making world lines...^^J}
  \foreach \i [evaluate={\x=\i*\d;}] in {1,...,\Nlines}{
    \message{  Running i/N=\i/\Nlines, x=\x...^^J}
    \draw[world line]   ( \x,-\ymin) -- ( \x,\ymax);
    \draw[world line t] (-\xmin, \x) -- (\xmax, \x);
  }
  \draw[world line]  (\dextra,-\ymin)  -- (\dextra,\ymax);
  \draw[world line]  (\dextra+\d,-\ymin) -- (\dextra+\d,\ymax);
  \draw[world line'] (-\xmin,\dextra) -- (\xmax,\dextra);
  
  % BOOSTED WORLD LINE GRID
  % \message{  Making world lines for boosted frame...^^J}
  % \fill[mydarkred,opacity=0.05]
  %   (O) --++ (\ang:\xmaxp) --++ (90-\ang:\xmaxp) --++ (\ang:-\xmaxp) -- cycle;
  % \fill[mydarkred,opacity=0.05]
  %   (O) --++ (\ang-180:\xminp) --++ (-90-\ang:\xminp) --++ (\ang:\xminp) -- cycle;
  % \foreach \i [evaluate={\x=\i*\D;}] in {1,...,\Nlines}{
  %   \message{  Running i/N=\i/\Nlines, x=\x...^^J};
  %   \draw[world line'] (\ang:\x) --++ (90-\ang:\xmaxp);
  %   \draw[world line'] (90-\ang:\x) --++ (\ang:\xmaxp);
  % }
  
  % AXES
  \draw[->,thick] (0,-\ymin) -- (T) node[above] {$ct$};
  \draw[->,thick] (-\xmin,0) -- (X) node[right] {$x$};
  \draw[->,thick,mydarkred] (-90-\ang:\xminp) -- (T')
    node[right=2,above=-1] {$ct'$};
  \draw[->,thick,mydarkred] (\ang-180:\xminp) -- (X') node[right] {$x'$};
  
  % % LIGHTCONE
  % \draw[myorange,thick]
  %   (-1.1*\xminp,1.1*\xminp) -- (1.1*\xminp,-1.1*\xminp)
  %   (-1.3*\xminp,-1.3*\xminp) -- (1.2*\xmaxp,1.2*\xmaxp);
  
  % % SPACELIKE HYPERBOLOIDS
  % \draw[mygreen,thick,samples=\Nsamples,smooth,variable=\x,domain=-1.6*\xmin:1.05*\xmax]
  %   plot(\x,{sqrt((\st)^2+(\x)^2)});
  % \draw[mydarkgreen,very thick,samples=\Nsamples,variable=\x,domain=0:\Ax,
  %       decoration={markings,mark=at position 0.58 with {\arrow{latex}}},postaction={decorate}]
  %   plot(\x,{sqrt((\st)^2+(\x)^2)});
  % \node[mydarkgreen,above left=-2] at (\xmax,{sqrt((\st)^2+(\xmax)^2)})
  %   {$s^2 = c^2t^2-x^2>0$};
  
  % % TIMELIKE HYPERBOLOIDS
  % \draw[myred,very thick,samples=\Nsamples,variable=\y,domain=\st:\Cx,
  %       decoration={markings,mark=at position 0.58 with {\arrow{latex}}},postaction={decorate}]
  %   plot({sqrt(\sr^2+(\y)^2)},\y);
  % \draw[myblue,thick,samples=\Nsamples,smooth,variable=\y,domain=-1.6*\ymin:0.95*\ymax]
  %   plot({sqrt(\sx^2+(\y)^2)},\y);
  % \draw[mydarkblue,very thick,samples=\Nsamples,variable=\y,domain=0:\By,
  %       decoration={markings,mark=at position 0.58 with {\arrow{latex}}},postaction={decorate}]
  %   plot({sqrt(\sx^2+(\y)^2)},\y);
  % \node[mydarkblue,below=0] at (0.7*\xmax,-1.6*\xmin)
  %   {$s^2 = c^2t^2-x^2<0$};
  
  % TICKS
  % \draw[mydarkgreen,dashed] ({\Ax},0) -- (A') -- (0,{\Ay});
  % \draw[mydarkblue,dashed] ({\Bx},0) -- (B') -- (0,{\By});
  % \tick{0,\st}{0} node[mydarkgreen,right=4,below left=-2] {$ct_1$};
  % \tick{\sx,0}{90} node[mydarkblue,below=1,below left=-3] {$x_1$};
  
  % EVENTS
  % \fill[mydarkgreen] (A)  circle(0.03); % event A
  % \fill[mydarkgreen] (A') circle(0.03); % event A'
  % \fill[mydarkblue]  (B)  circle(0.03); % event B
  % \fill[mydarkblue]  (B') circle(0.03); % event B'
  % \fill[mydarkred]   (C)  circle(0.03); % event C
  % \fill[mydarkred] (\ang:4*\D)++(90-\ang:3*\D) coordinate (C') circle(0.03); % event C'
}

% SPACETIME DIAGRAM - SIMULTANEITY (in S)
\begin{tikzpicture}[scale=1.8]
  \message{Simultaneity^^J}
  
  % AXES
  
  % SETTINGS
  \axes
  \def\L{0.91*\xmaxp} % length of the dashed lines
  \def\fact{3}
  \pgfmathsetmacro\xA{2*\d} % x coordinate of A in S
  \pgfmathsetmacro\yA{3*\d} % x coordinate of A in S
  \pgfmathsetmacro\dx{\fact*\d} % time difference in S
  \pgfmathsetmacro\xAp{(\xA-tan(\ang)*\yA)/cos(\ang)^2/sqrt(1-tan(\ang)^2)} % x coordinate of A in S'
  \pgfmathsetmacro\dtp{\fact*\d*sin(\ang)/cos(2*\ang)} % time difference between A and B in S'
  \coordinate (A) at (\xA,\yA);
  \coordinate (B) at (\xA+\dx,\yA);
  \coordinate (A') at ($(A)-(\ang:0.13*\xmaxp)$); % left side of dashed line through A
  \coordinate (B') at ($(A')-(90-\ang:\dtp)$); % left side of dashed line through B
  
  % FILL
  % \fill[mylightred]
  %   (A)++(\ang:-\xAp) --++ (90-\ang:-\dtp) -- ($(B)+(\ang:0.07)$) --++ (90-\ang:\dtp) -- cycle;
  % \draw[mylightblue2,line width=1.8] (0,\yA) -- (B);
  
  % EVENTS
  %\draw[mygreen,thick] (A) -- (B);
  \draw[mydarkgreen,mydashed,thin]
  (A) -- (A') -- ++ (90+2*\ang:0.1) node[left,fill=white,inner sep=0.1em] {$ct'_1$};
  \draw[mydarkblue,mydashed,thin]
  (B) -- (B') -- ++ (90+2*\ang:0.1) node[left,fill=white,inner sep=0.1em] {$ct'_2$};
  \draw[mydarkgreen,mydashed,thin]
  (A) -- (\xA,0) node[below] {\contour{white}{$x_1$}};
  \draw[mydarkblue,mydashed,thin]
  (B) -- (\xA+\dx,0) node[below] {\contour{white}{$x_2$}};
  \draw[dashed] (B) -- (0,\yA) node[left=2,fill=white,inner sep=0.1em] {$ct_1=ct_2$};
  \draw[mydarkorange,thick,<->] (A) -- (B)
  node[pos=0.5,above] {$\Delta x$};
  \fill[mydarkgreen] (A) circle(0.04) % event A
  node[above=0] {\contour{white}{$E_1$}};
  \fill[mydarkblue] (B) circle(0.04) % event B
  node[above] {\contour{white}{$E_2$}};
  
  % ARROW
  \draw[<->,mydarkorange,thick] (A') -- (B')
  node[pos=0.3,left=3,fill=white,inner sep =0em] {\contour{white}{$c\Delta t'$}};
  
\end{tikzpicture}
\end{document}